\uuid{gqgj}
\titre{Estimation d'un paramètre de la loi de Weibull}
\niveau{L3}
\module{Probabilité et statistique}
\chapitre{Estimateur}
\sousChapitre{Estimateur}
\theme{Loi de Weibull, maximum de vraisemblance, loi du chi-deux, biais, convergence}
\auteur{Erwan L'haridon}
\datecreate{2026-05-19}
\organisation{AMSCC}
\difficulte{3}

\contenu{
	
	\texte{
		La durée de vie d'un matériel est modélisée par une variable aléatoire $X$ suivant une loi de Weibull $\mathcal{W}\!\left(\lambda, \frac{1}{2}\right)$ de densité :
		$$
		f(x, \lambda) = \frac{1}{2\lambda\sqrt{x}}\exp\!\left(-\frac{\sqrt{x}}{\lambda}\right)\mathbf{1}_{\mathbb{R}^+}(x),
		$$
		dont l'espérance vaut $\mathbb{E}(X) = 2\lambda$ et la variance $\mathbb{V}(X) = 20\lambda^2$.
		
		Soit $(X_1, \ldots, X_n)$ un $n$-échantillon i.i.d. de loi $\mathcal{W}\!\left(\lambda, \frac{1}{2}\right)$.
	}
	
	\begin{enumerate}
		
		\item
		\question{Déterminer l'estimateur du maximum de vraisemblance $\hat{\lambda}$ de $\lambda$.}
		\indication{Étudier la log-vraisemblance et annuler sa dérivée par rapport à $\lambda$.}
		\reponse{
			Pour $x_i > 0$, $\forall\, 1 \leq i \leq n$, la vraisemblance est :
			\[
			\mathcal{L}(x_1, \ldots, x_n, \lambda)
			= \prod_{i=1}^n f(x_i, \lambda)
			= \frac{1}{2^n \lambda^n \sqrt{\prod_{i=1}^n x_i}}
			\exp\!\left(-\frac{1}{\lambda}\sum_{i=1}^n \sqrt{x_i}\right).
			\]
			La log-vraisemblance est :
			\[
			\hat{\mathcal{L}}(x_1, \ldots, x_n, \lambda)
			= -n\ln 2 - n\ln\lambda - \frac{1}{2}\sum_{i=1}^n \ln x_i - \frac{1}{\lambda}\sum_{i=1}^n \sqrt{x_i}.
			\]
			On annule la dérivée par rapport à $\lambda$ :
			\[
			\frac{\partial}{\partial\lambda}\hat{\mathcal{L}} = -\frac{n}{\lambda} + \frac{1}{\lambda^2}\sum_{i=1}^n \sqrt{x_i} = 0
			\iff \lambda = \frac{1}{n}\sum_{i=1}^n \sqrt{x_i}.
			\]
			On vérifie qu'il s'agit bien d'un maximum en calculant la dérivée seconde évaluée en $\hat{\lambda}$ :
			\[
			\frac{\partial^2}{\partial\lambda^2}\hat{\mathcal{L}}\!\left(x_1,\ldots,x_n,\frac{1}{n}\sum_{i=1}^n\sqrt{x_i}\right)
			= -\frac{n^3}{\left(\sum_{i=1}^n\sqrt{x_i}\right)^2} < 0.
			\]
			L'estimateur du maximum de vraisemblance est donc :
			\[
			\hat{\lambda} = \frac{1}{n}\sum_{i=1}^n \sqrt{X_i}.
			\]
		}
		
		\item
		\question{Montrer que $\dfrac{2}{\lambda}\sqrt{X}$ suit une loi du $\chi^2$ dont on précisera le paramètre. En déduire que $\dfrac{2}{\lambda}\displaystyle\sum_{i=1}^n \sqrt{X_i} \sim \chi^2(2n)$.}
		\indication{Pour toute fonction $\varphi$ continue bornée, calculer $\mathbb{E}\!\left(\varphi\!\left(\frac{2}{\lambda}\sqrt{X}\right)\right)$ à l'aide du changement de variable $u = \frac{2}{\lambda}\sqrt{x}$.}
		\reponse{
			Soit $\varphi$ une fonction continue bornée. On calcule :
			\begin{align*}
				\mathbb{E}\!\left(\varphi\!\!\left(\frac{2}{\lambda}\sqrt{X}\right)\right)
				&= \int_0^{+\infty} \varphi\!\!\left(\frac{2}{\lambda}\sqrt{x}\right)
				\frac{1}{2\lambda\sqrt{x}}\exp\!\left(-\frac{\sqrt{x}}{\lambda}\right)dx.
			\end{align*}
			On pose $u = \dfrac{2}{\lambda}\sqrt{x}$, soit $x = \dfrac{\lambda^2 u^2}{4}$ et $dx = \dfrac{\lambda^2 u}{2}\,du$. Alors $\sqrt{x} = \dfrac{\lambda u}{2}$ et :
			\begin{align*}
				\mathbb{E}\!\left(\varphi\!\!\left(\frac{2}{\lambda}\sqrt{X}\right)\right)
				&= \int_0^{+\infty} \varphi(u)\,
				\frac{1}{2\lambda \cdot \frac{\lambda u}{2}}
				\exp\!\left(-\frac{u}{2}\right)
				\frac{\lambda^2 u}{2}\,du \\
				&= \int_0^{+\infty} \varphi(u)\,\frac{1}{2}\exp\!\left(-\frac{u}{2}\right)du \\
				&= \int_0^{+\infty} \varphi(u)\,\frac{1}{\Gamma\!\left(\frac{1}{2}\right) \cdot 2^{1/2}}\,
				u^{\frac{1}{2}-1} e^{-u/2}\,du \cdot \Gamma\!\!\left(\tfrac{1}{2}\right)2^{1/2} \cdot \frac{1}{2}.
			\end{align*}
			On reconnaît la densité d'une loi $\chi^2(2)$ (i.e. $\Gamma\!\left(1, \frac{1}{2}\right)$) :
			\[
			\frac{1}{2}e^{-u/2} = \frac{1}{\Gamma(1)\cdot 2^1}\,u^{1-1}e^{-u/2}.
			\]
			Donc $\dfrac{2}{\lambda}\sqrt{X} \sim \chi^2(2)$.
			
			Les variables $\dfrac{2}{\lambda}\sqrt{X_i}$ sont i.i.d. de loi $\chi^2(2)$. Par stabilité de la loi du $\chi^2$ par sommation de variables indépendantes :
			\[
			\frac{2}{\lambda}\sum_{i=1}^n \sqrt{X_i} \sim \chi^2(2n).
			\]
		}
		
		\item
		\question{En déduire la loi de $\dfrac{2n\hat{\lambda}}{\lambda}$, puis le biais et la variance de $\hat{\lambda}$. L'estimateur $\hat{\lambda}$ est-il convergent ?}
		\reponse{
			On a $\dfrac{2n\hat{\lambda}}{\lambda} = \dfrac{2}{\lambda}\displaystyle\sum_{i=1}^n \sqrt{X_i} \sim \chi^2(2n)$.
			
			Si $T \sim \chi^2(m)$, alors $\mathbb{E}(T) = m$ et $\mathbb{V}(T) = 2m$. On en déduit :
			\[
			\mathbb{E}\!\left(\frac{2n\hat{\lambda}}{\lambda}\right) = 2n
			\implies \frac{2n}{\lambda}\,\mathbb{E}(\hat{\lambda}) = 2n
			\implies \mathbb{E}(\hat{\lambda}) = \lambda.
			\]
			Donc $\hat{\lambda}$ est \textbf{sans biais} : $\mathbb{B}(\hat{\lambda}) = 0$.
			
			Pour la variance :
			\[
			\mathbb{V}\!\left(\frac{2n\hat{\lambda}}{\lambda}\right) = 4n
			\implies \frac{4n^2}{\lambda^2}\,\mathbb{V}(\hat{\lambda}) = 4n
			\implies \mathbb{V}(\hat{\lambda}) = \frac{\lambda^2}{n}.
			\]
			L'estimateur étant sans biais, le risque quadratique est :
			\[
			\mathrm{EQM}(\hat{\lambda}) = \mathbb{V}(\hat{\lambda}) = \frac{\lambda^2}{n} \xrightarrow[n\to+\infty]{} 0.
			\]
			Donc $\hat{\lambda}$ est \textbf{convergent} (en moyenne quadratique, donc en probabilité).
		}
		
	\end{enumerate}
}